\documentclass[../../../../main.tex]{subfiles}
\begin{document}

\begin{flushright}
\footnotesize\textit{【自動文字起こし・要確認】}
\end{flushright}

{\bf [解]}
\textbf{(1)} $O$を原点とし，$\vec a=\overrightarrow{OA}$，$\vec b=\overrightarrow{OB}$，$\vec c=\overrightarrow{OC}$（$|\vec a|=|\vec b|=|\vec c|=1$，$\vec a\cdot\vec b=\vec b\cdot\vec c=\vec c\cdot\vec a=\dfrac12$）とおく．$D$は$AB$の中点，$E$は$OC$の中点だから
\begin{align*}
\overrightarrow{AC}=\vec c-\vec a,\qquad\overrightarrow{DE}=\frac{\vec c}2-\frac{\vec a+\vec b}2=\frac12(\vec c-\vec a-\vec b).
\end{align*}
よって
\begin{align*}
\overrightarrow{AC}\cdot\overrightarrow{DE}&=\frac12\bigl\{(\vec c-\vec a)\cdot(\vec c-\vec a)-(\vec c-\vec a)\cdot\vec b\bigr\}=\frac12\bigl\{|\vec c-\vec a|^2-\vec b\cdot\vec c+\vec a\cdot\vec b\bigr\}.
\end{align*}
$|\vec c-\vec a|^2=|\vec c|^2-2\vec a\cdot\vec c+|\vec a|^2=1-1+1=1$，$\vec b\cdot\vec c=\vec a\cdot\vec b=\dfrac12$だから
\begin{align*}
\overrightarrow{AC}\cdot\overrightarrow{DE}=\frac12\Bigl(1-\frac12+\frac12\Bigr)=\frac12.
\end{align*}

\textbf{(2)} 目の積が10の倍数（$2$の倍数かつ$5$の倍数）にならない事象を考える．$B$:「3回とも偶数の目が出ない（目が$1,3,5$のみ）」，$C$:「3回とも$5$が出ない（目が$1,2,3,4,6$のみ）」，$A=B\cap C$:「3回とも目が$1,3$のみ」とすると，求める事象の余事象は$B\cup C$である．
\begin{align*}
P(A)=\Bigl(\frac26\Bigr)^3=\Bigl(\frac13\Bigr)^3,\qquad P(B)=\Bigl(\frac36\Bigr)^3=\Bigl(\frac12\Bigr)^3,\qquad P(C)=\Bigl(\frac56\Bigr)^3.
\end{align*}
包除原理より$P(B\cup C)=P(B)+P(C)-P(B\cap C)=P(B)+P(C)-P(A)$だから，求める確率は
\begin{align*}
P(\overline{B\cup C})=1-P(B)-P(C)+P(A)=1-\Bigl(\frac12\Bigr)^3-\Bigl(\frac56\Bigr)^3+\Bigl(\frac13\Bigr)^3=1-\frac18-\frac{125}{216}+\frac1{27}=\frac{216-27-125+8}{216}=\frac{72}{216}=\frac13.
\end{align*}

（別解：直接数え上げる．3個の目の中に「5の倍数」（=5）と「2の倍数」（=2,4,6）が少なくとも1つずつ含まれる場合を数える．5の倍数の個数と2の倍数の個数の内訳で場合分けすると，$3!\cdot1\cdot3\cdot2=36$，${}_3C_1\cdot3\cdot1=9$，${}_3C_1\cdot1\cdot3=27$の3通りの内訳があり，合計$36+9+27=72$通り．よって確率は$72/6^3=72/216=1/3$．）


\newpage

\end{document}
